
\medskip
\phantomsection\label{prop:sga3-VIB-7.8}%
\noindent\textbf{Proposition 7.8.---}
\emph{Let \(G\) be a \(k\)-group locally of finite type, and
let \(A\) and \(B\) be two geometrically reduced \(k\)-group
subschemes of \(G\) (hence smooth at the generic points of
their irreducible components, and therefore smooth by
\(1.3\)). Suppose that one of the following conditions holds:}
\begin{enumerate}
\item[\textup{a)}]\phantomsection\label{cond:sga3-VIB-7.8-a}
  \emph{\(A\) and \(B\) are invariant and of finite type over
  \(k\);}
\item[\textup{b)}]\phantomsection\label{cond:sga3-VIB-7.8-b}
  \emph{\(A\) is connected and \(B\) is of finite type over
  \(k\).}
\end{enumerate}
\emph{Then \((A,B)\) is of finite type over \(k\), and
represents the sheaf associated, for the \((\mathrm{fppf})\)
topology (or the \((\mathrm{fpqc})\) topology), with the group
presheaf of commutators of \(A\) and \(B\) in \(G\).
Moreover,}%
\footnote{\label{note:sga3-VIB-79}Editor's note: the assertions
that follow have been made explicit, and the reduction to the
algebraically closed case has been expanded in the proof.}
\emph{the \(k\)-groups \((A,B^0)\) and \((A^0,B)\) are
connected, and}
\phantomsection\label{form:sga3-VIB-7.8-identity-component}%
\[
  (A,B)^0=(A,B^0)\cdot(A^0,B).
\]

\emph{Proof.}
By \hyperref[prop:sga3-VIB-7.6]{7.6}, to show that \((A,B)\)
is the desired associated sheaf, it is enough to show that
there exists an integer \(n\) such that
\phantomsection\label{form:sga3-VIB-7.8-finite-generation}%
\[
  \nu^n\bigl((A\times_kB)^n\bigr)=(A,B),
\]
with the notation of
\hyperref[def:sga3-VIB-7.2.2]{7.2.2}. To prove this and the
other two assertions, we may suppose that \(k\) is
algebraically closed.

Indeed, let \(\overline{k}\) be an algebraic closure of \(k\).
By \hyperref[concl:sga3-VIB-7.1-iii]{7.1\textup{(iii)}} and
\hyperref[prop:sga3-VIA-2.4]{VI\(_A\), 2.4}, one has, with
the evident notation,
\phantomsection\label{form:sga3-VIB-7.8-base-change}%
\[
\begin{gathered}
  (A,B)_{\overline{k}}
    =(A_{\overline{k}},B_{\overline{k}}),\\
  \bigl((A,B)^0\bigr)_{\overline{k}}
    =(A_{\overline{k}},B_{\overline{k}})^0,\\
  (A,B^0)_{\overline{k}}
    =(A_{\overline{k}},B_{\overline{k}}^0),
  \qquad\text{and so on.}
\end{gathered}
\]
Consequently, suppose that
\((A_{\overline{k}},B_{\overline{k}})\) is of finite type over
\(\overline{k}\); respectively, suppose that
\((A_{\overline{k}},B_{\overline{k}}^0)\) and
\((A_{\overline{k}}^0,B_{\overline{k}})\) are connected and
that the morphism\begingroup
\renewcommand{\thefootnote}{\fnsymbol{footnote}}%
\footnote[1]{\label{note:sga3-VIB-7.8-mixed-base-change}
Translator's note: the French source prints
\((A_{\overline{k}},B_{\overline{k}}^0)\) in both factors of
this product. The second factor must be
\((A_{\overline{k}}^0,B_{\overline{k}})\), as in the asserted
identity and throughout the proof; the mixed product is
restored here.}%
\endgroup
\[
  (A_{\overline{k}},B_{\overline{k}}^0)
  \times_{\overline{k}}
  (A_{\overline{k}}^0,B_{\overline{k}})
  \longrightarrow
  (A_{\overline{k}},B_{\overline{k}})^0
\]
is surjective. The analogous assertions then hold over \(k\),
by EGA~IV\(_2\), 2.7.1 and 2.6.1.

Let
\[
  B^1,\ldots,B^p
\]
be the connected components of \(B\) other than the identity
component \(B^0\). There are only finitely many, since \(B\)
is of finite type over \(k\), hence Noetherian. In case
\textup{a)}, likewise let
\[
  A^1,\ldots,A^q
\]
be the connected components of \(A\) other than \(A^0\). In
case \textup{b)}, only \(A^0=A\) is considered. Let
\(\nu_{ij}\) be the restriction of \(\nu\) to
\[
  A^i\times_kB^j.
\]
Each \(A^i\) and \(B^j\) is irreducible, by
\hyperref[cor:sga3-VIA-2.4.1]{VI\(_A\), 2.4.1}; hence so are
\(A^0\times_kB^j\) and \(A^i\times_kB^0\). Since the unit
element of \(G\) belongs to \(A^0\) and to \(B^0\), it belongs
to
\[
  \nu_{0j}(A^0\times_kB^j)
  \qquad\text{and}\qquad
  \nu_{i0}(A^i\times_kB^0).
\]
Thus every \(\Gamma_G(\nu_{0j})\) and every
\(\Gamma_G(\nu_{i0})\) is connected, by
\hyperref[cor:sga3-VIB-7.2.1]{7.2.1}.

Let \(u_{0j}\), respectively \(u_{i0}\), denote the inclusion
of \(\Gamma_G(\nu_{0j})\), respectively
\(\Gamma_G(\nu_{i0})\), in \(G\). Then
\phantomsection\label{form:sga3-VIB-7.8-generated-connected-groups}%
\[
  (A^0,B)=\Gamma_G\bigl((u_{0j})_{j=0}^{p}\bigr),
  \qquad
  (A,B^0)=\Gamma_G\bigl((u_{i0})_{i=0}^{q}\bigr),
\]
and these groups are connected. Moreover, it follows from
\hyperref[prop:sga3-VIB-7.4]{7.4} and the preceding
constructions that there is an integer \(r\) such that both
\((A^0,B)\) and \((A,B^0)\) are contained in
\[
  \nu^r\bigl((A\times_kB)^r\bigr).
\]
In case \textup{b)}, one has \((A,B)=(A^0,B)\), and the proof
is complete.

We now consider case \textup{a)}.%
\footnote{\label{note:sga3-VIB-80}Editor's note: the argument
that follows has been expanded to take account of the added
\hyperref[cor:sga3-VIB-7.1.1]{Corollary 7.1.1}.}
We already know that \((A^0,B)\) and \((A,B^0)\) are smooth
connected \(k\)-subgroups of \(G\), hence are of finite type
over \(k\)
(\hyperref[prop:sga3-VIA-2.4]{VI\(_A\), 2.4}). Since \(A^0\)
is a characteristic subgroup of \(A\)
(\hyperref[thm:sga3-VIA-2.6.5]{VI\(_A\), 2.6.5}), it is an
invariant subgroup of \(G\); hence
\(\hyperref[concl:sga3-VIB-7.3-v]{7.3\textup{(v)}}\) shows
that \((A^0,B)\) is an invariant subgroup of \(G\), and the
same holds for \((A,B^0)\).

By \hyperref[cor:sga3-VIB-7.1.1]{7.1.1}, the subgroup \(H\) of
\(G\) generated by \((A,B^0)\) and \((A^0,B)\) is therefore
\[
  H=(A,B^0)\cdot(A^0,B).
\]
In particular,
\phantomsection\label{form:sga3-VIB-7.8-H-contained}%
\[
  H\subset\nu^{2r}\bigl((A\times_kB)^{2r}\bigr).
\]

If \(T\) is a subset of \(G\) stable under the group law
(cf.\ \(3.0\)), write \(T_0\) for the group of \(k\)-points of
\(G\) belonging to \(T\). Put
\[
  \Gamma'
    =\bigcup_{q\geq1}\nu^q\bigl((A\times_kB)^q\bigr).
\]
By \hyperref[prop:sga3-VIB-7.9]{Proposition 7.9} below,
\phantomsection\label{form:sga3-VIB-7.8-rational-commutators}%
\[
  (A^0,B)_0=(A_0^0,B_0),
  \qquad
  (A,B^0)_0=(A_0,B_0^0),
  \qquad
  \Gamma'_0=(A_0,B_0).
\]
Therefore
\[
  H_0=(A_0^0,B_0)\cdot(A_0,B_0^0)
\]
has finite index in \(\Gamma'_0\) (Bible, Expos\'e~III,
Appendix), since \(A_0\) and \(B_0\) are invariant and
\(A_0^0\), respectively \(B_0^0\), has finite index in
\(A_0\), respectively \(B_0\).

The hypotheses of \hyperref[lem:sga3-VIB-7.7]{Lemma 7.7} now
hold. Since
\[
  H\subset\nu^{2r}\bigl((A\times_kB)^{2r}\bigr),
\]
there exists an integer \(m\) such that
\[
  \nu^{2rm}\bigl((A\times_kB)^{2rm}\bigr)=(A,B),
\]
and there are finitely many \(k\)-points
\(a_1,\ldots,a_n\) of \(G\) such that
\[
  (A,B)=a_1H\cup\cdots\cup a_nH.
\]
Since \(H\) is of finite type over \(k\), every \(a_iH\) is
quasi-compact. Their union \((A,B)\) is therefore
quasi-compact, hence of finite type over \(k\). Finally, \(H\)
is irreducible, so each \(a_iH\) is irreducible; and because
the unit \(e\) belongs to \(H\), it follows that
\[
  H=(A,B)^0.
\]
